% =============================================================================
%  CHAPTER 3 -- LINEAR REGRESSION FROM SCRATCH
% =============================================================================
\section{Linear regression from scratch}
\label{ch:ols}

Regression is the only estimation tool this project needs, and it is used in four
different places: to build the hedge, to test for unit roots, to fit the
Ornstein--Uhlenbeck process, and to run the Kalman filter's update step. We therefore
derive it completely.

\subsection{The setup}
\label{sec:ols:setup}

We observe $T$ pairs $(y_t, x_t)$ with $y_t\in\R$ and $x_t\in\R^{n}$, and posit
\begin{equation}
  y_t \;=\; x_t^{\top}\beta + \eps_t,\qquad t=1,\dots,T,
  \label{eq:ols:model}
\end{equation}
where $\beta\in\R^n$ is unknown and $\eps_t$ is an unobserved error. Stack the data:
\begin{equation}
  Y=\begin{bmatrix}y_1\\ \vdots\\ y_T\end{bmatrix}\in\R^{T},\qquad
  X=\begin{bmatrix}x_1^{\top}\\ \vdots\\ x_T^{\top}\end{bmatrix}\in\R^{T\times n},
  \qquad
  \eps=\begin{bmatrix}\eps_1\\ \vdots\\ \eps_T\end{bmatrix}\in\R^{T},
\end{equation}
so that
\begin{equation}
  Y = X\beta + \eps,\qquad \eps\sim(0,\ \sigma^{2}I_T).
  \label{eq:linmodel}
\end{equation}
The notation $\eps\sim(0,\sigma^2 I)$ means $\E[\eps]=0$ and
$\Var(\eps)=\E[\eps\eps^{\top}]=\sigma^2 I$: errors are homoskedastic
(same variance $\sigma^2$) and serially uncorrelated (off-diagonal zeros). If we want an
intercept, we include a column of ones in $X$ and let the first element of $\beta$ be
$\alpha$; nothing else changes.

In our application: $y_t$ is the target's log price, $x_t=[1,\ \log P^{1}_t,\dots,\log
P^{n}_t]^{\top}$, and $\beta=[\alpha,\beta_1,\dots,\beta_n]^{\top}$. The residual
$\hat\eps_t$ is the spread of equation~\eqref{eq:saa:spread}.

\subsection{Least squares as geometry}
\label{sec:ols:geometry}

\emph{Column space.} The set of all vectors that can be written as $X\beta$ for some
$\beta$ is the column space $\mathcal C(X)=\operatorname{span}\{X_{\cdot 1},\dots,X_{\cdot n}\}$,
an $n$-dimensional subspace of $\R^T$ (assuming full column rank). Choosing $\beta$ is
choosing a point in that subspace; the fitted values $\hat Y=X\hat\beta$ are that point.

\emph{The objective.} Least squares picks the point of $\mathcal C(X)$ closest to $Y$ in
Euclidean distance:
\begin{equation}
  \hat\beta \;=\; \argmin_{\beta}\ \lVert Y-X\beta\rVert_2^2
  \;=\; \argmin_{\beta}\ (Y-X\beta)^{\top}(Y-X\beta).
  \label{eq:ols:objective}
\end{equation}

\begin{intuition}
The closest point in a subspace is the \emph{orthogonal projection}: the residual
$\hat\eps=Y-X\hat\beta$ must be perpendicular to the subspace, i.e.\ perpendicular to
every column of $X$. Perpendicular means inner product zero: $X^{\top}\hat\eps=0$. That
single geometric requirement \emph{is} the normal equations, and it is worth internalising
because it explains half of regression's properties: residuals sum to zero (they are
orthogonal to the ones-column), residuals are uncorrelated with the regressors
\emph{in sample}, and $R^2$ decomposes exactly the way it does.
\end{intuition}

\subsection{Two derivations of the normal equations}
\label{sec:ols:normal}

\paragraph{(a) Calculus.} Expand the objective:
\[
  Q(\beta) = Y^{\top}Y - 2\beta^{\top}X^{\top}Y + \beta^{\top}X^{\top}X\beta .
\]
Differentiate using \eqref{eq:matrixcalculus} ($X^{\top}X$ is symmetric):
\[
  \nabla_\beta Q = -2X^{\top}Y + 2X^{\top}X\beta .
\]
Set to zero:
\begin{equation}
  X^{\top}X\,\hat\beta \;=\; X^{\top}Y
  \qquad\text{(the normal equations)}.
  \label{eq:normal}
\end{equation}
The Hessian is $\nabla^2_\beta Q = 2X^{\top}X$, which is positive semi-definite for any
$X$ (since $a^{\top}X^{\top}Xa=\lVert Xa\rVert^2\ge0$), so $Q$ is convex and any
stationary point is a global minimum.

\paragraph{(b) Geometry.} Impose $X^{\top}(Y-X\hat\beta)=0$. This is literally
\eqref{eq:normal}. No calculus needed.

\paragraph{Closed form.} If $X^{\top}X$ is invertible (equivalently $\rank(X)=n$,
equivalently no exact linear dependence among the columns),
\begin{equation}
  \boxed{\ \hat\beta \;=\; \bigl(X^{\top}X\bigr)^{-1}X^{\top}Y\ }
  \label{eq:ols}
\end{equation}
and the minimizer is unique. If $X^{\top}X$ is singular, \eqref{eq:normal} still has
solutions (since $X^{\top}Y\in\mathcal C(X^{\top}X)$) but they are not unique: any two
solutions give the same fitted values $\hat Y$, and the minimum-norm one is given by the
Moore--Penrose pseudo-inverse $X^{+}$. In practice we never rely on that: the ridge
penalty of Section~\ref{sec:ridge} adds $\lambda I$ and restores invertibility, which is
also why ridge fixes ``singular $X^{\top}X$'' crashes for free.

\paragraph{The hat matrix.} Define
\begin{equation}
  H \;=\; X\bigl(X^{\top}X\bigr)^{-1}X^{\top}\in\R^{T\times T},
  \qquad \hat Y = HY,\qquad \hat\eps = (I-H)Y .
  \label{eq:hatmatrix}
\end{equation}
$H$ is the orthogonal projector onto $\mathcal C(X)$. Its properties, each of which is
used somewhere in this book:
\begin{align}
  H^2 &= H \quad\text{(idempotent: projecting twice is projecting once)}, \label{eq:hat:idem}\\
  H^{\top} &= H \quad\text{(symmetric)}, \label{eq:hat:sym}\\
  \tr(H) &= \tr\bigl((X^{\top}X)^{-1}X^{\top}X\bigr) = n, \label{eq:hat:trace}\\
  \tr(I-H) &= T-n, \label{eq:hat:trace2}\\
  X^{\top}(I-H) &= 0 \quad\text{(residuals orthogonal to regressors)}. \label{eq:hat:orth}
\end{align}
The diagonal $h_{tt}=x_t^{\top}(X^{\top}X)^{-1}x_t$ is the \emph{leverage} of
observation $t$: how much that observation influences its own fitted value. It satisfies
$0\le h_{tt}\le1$ and $\sum_t h_{tt}=n$, so the average leverage is $n/T$. High-leverage
points are observations far from the regressor cloud; in a stat-arb context they are the
days when the basket configuration was unusual, and they are exactly where a single bad
print can distort the hedge.

\subsection{Statistical properties: unbiasedness and variance}
\label{sec:ols:properties}

\begin{theorem}[Unbiasedness]
If $\E[\eps\mid X]=0$ and $\rank(X)=n$, then $\E[\hat\beta\mid X]=\beta$.
\end{theorem}
\begin{proof}
Substitute $Y=X\beta+\eps$ into \eqref{eq:ols}:
\[
  \hat\beta = (X^{\top}X)^{-1}X^{\top}(X\beta+\eps)
  = \beta + (X^{\top}X)^{-1}X^{\top}\eps .
\]
Take conditional expectations: $\E[\hat\beta\mid X]=\beta+(X^{\top}X)^{-1}X^{\top}\E[\eps\mid X]=\beta$.
\end{proof}

The second term, $A\eps$ with $A=(X^{\top}X)^{-1}X^{\top}$, is the \emph{estimation
noise}. Its conditional covariance is
\begin{equation}
  \Var(\hat\beta\mid X) \;=\; \E\bigl[A\eps\eps^{\top}A^{\top}\mid X\bigr]
  \;=\; A\,\E[\eps\eps^{\top}\mid X]\,A^{\top}
  \;\overset{\Var(\eps)=\sigma^2 I}{=}\; \sigma^{2}AA^{\top}
  \;=\; \sigma^{2}\bigl(X^{\top}X\bigr)^{-1},
  \label{eq:ols:varbeta}
\end{equation}
where the last step uses $AA^{\top}=(X^{\top}X)^{-1}X^{\top}X(X^{\top}X)^{-1}=(X^{\top}X)^{-1}$.

Write $X^{\top}X = \sum_{t} x_tx_t^{\top} = T\hat\Sigma_{xx}$ where $\hat\Sigma_{xx}$
is the sample second-moment matrix of the regressors. Then
\begin{equation}
  \Var(\hat\beta_j) \;=\; \frac{\sigma^{2}}{T}\,\bigl[\hat\Sigma_{xx}^{-1}\bigr]_{jj},
  \label{eq:ols:varscale}
\end{equation}
which displays the three things that make a coefficient uncertain: the noise level
$\sigma^2$, the sample size $T$ (as $1/T$ in variance, $1/\sqrt T$ in standard error ---
the $\sqrt T$ of Section~\ref{sec:prob:clt} again), and the regressor geometry through
$\hat\Sigma_{xx}^{-1}$, which blows up when regressors are nearly collinear.

\begin{theorem}[Gauss--Markov]
\label{thm:gaussmarkov}
Under the classical assumptions
\begin{enumerate}[label=(G\arabic*),leftmargin=*]
  \item linearity: $Y=X\beta+\eps$;
  \item strict exogeneity: $\E[\eps\mid X]=0$;
  \item spherical errors: $\Var(\eps\mid X)=\sigma^2 I_T$;
  \item full rank: $\rank(X)=n$,
\end{enumerate}
$\hat\beta$ is the \emph{best linear unbiased estimator} (BLUE): for any other estimator
$\tilde\beta=CY$ that is unbiased for all $\beta$, $\Var(\tilde\beta)-\Var(\hat\beta)$ is
positive semi-definite.
\end{theorem}
\begin{proof}[Proof sketch]
Unbiasedness of $CY$ for all $\beta$ requires $CX=I$. Write $C=(X^{\top}X)^{-1}X^{\top}+D$;
then $CX=I+DX=I$ forces $DX=0$. Hence
$\Var(CY)=\sigma^2 CC^{\top}=\sigma^2\bigl[(X^{\top}X)^{-1}+DD^{\top}\bigr]$
(the cross terms vanish because $D X (X^{\top}X)^{-1}X^{\top}=0$), and $DD^{\top}$ is PSD.
So $\Var(CY)\succeq\Var(\hat\beta)$ with equality iff $D=0$.
\end{proof}

\begin{remark}[What Gauss--Markov does \emph{not} say]
It says nothing about biased estimators --- ridge and lasso are not BLUE and can still be
better in MSE (equation~\eqref{eq:biasvariance}). It says nothing about the
\emph{distribution} of $\hat\beta$ (for that you need normality of $\eps$). And it is
silent about the case where (G3) fails --- which is the normal case in finance. When
$\Var(\eps\mid X)=\Omega\ne\sigma^2I$, the correct covariance is the \emph{sandwich}
\begin{equation}
  \Var(\hat\beta\mid X) \;=\; \bigl(X^{\top}X\bigr)^{-1}X^{\top}\Omega X\bigl(X^{\top}X\bigr)^{-1},
  \label{eq:sandwich}
\end{equation}
and the OLS estimator, while still unbiased, is no longer efficient --- GLS,
$(X^{\top}\Omega^{-1}X)^{-1}X^{\top}\Omega^{-1}Y$, is. Equation~\eqref{eq:sandwich} is
the skeleton of every robust-standard-error formula in this book, including
Newey--West \eqref{eq:hac} and the White heteroskedasticity test.
\end{remark}

\subsection{Estimating $\sigma^2$ and the degrees-of-freedom correction}
\label{sec:ols:sigma2}

The natural candidate is $\hat\eps^{\top}\hat\eps/T$, and it is biased. Derivation:
$\hat\eps=(I-H)\eps$, so
\[
  \E[\hat\eps^{\top}\hat\eps] = \E\bigl[\tr(\hat\eps\hat\eps^{\top})\bigr]
  = \E\bigl[\tr((I-H)\eps\eps^{\top}(I-H))\bigr]
  = \tr\bigl((I-H)\,\sigma^2 I\,(I-H)\bigr)
  = \sigma^2\tr\bigl((I-H)^2\bigr)
  = \sigma^2\tr(I-H) = \sigma^2(T-n),
\]
using cyclicity of the trace, idempotence \eqref{eq:hat:idem}, and
\eqref{eq:hat:trace2}. Therefore the unbiased estimator is
\begin{equation}
  \hat\sigma^{2} \;=\; \frac{\hat\eps^{\top}\hat\eps}{T-n}
  \;=\; \frac{\mathrm{RSS}}{T-n},
  \label{eq:sigma2hat}
\end{equation}
exactly the $T-1$ argument of Section~\ref{sec:prob:variance} generalised to $n$
estimated coefficients. \emph{In our backtests} the window is $W=120$ bars and $n=4$
(intercept plus three basket names), so the divisor is 116, not 120 --- a $3.4\%$
correction in $\hat\sigma^2$, $1.7\%$ in $\hat\sigma$. Small, but it is the difference
between the spread $z$-score you compute and the one you think you compute.

Under normality, $(T-n)\hat\sigma^2/\sigma^2\sim\chi^2_{T-n}$ independently of
$\hat\beta$, hence
\begin{equation}
  \frac{\hat\beta_j-\beta_j}{\hat s_j}\;\sim\; t_{T-n},
  \qquad \hat s_j = \sqrt{\hat\sigma^2\,\bigl[(X^{\top}X)^{-1}\bigr]_{jj}},
  \label{eq:tstat}
\end{equation}
which is the standard $t$-statistic reported for each hedge weight --- and the object
that HAC estimation replaces in Section~\ref{sec:ts:hac}.

\subsection{$R^2$, adjusted $R^2$, and the $F$-test}
\label{sec:ols:r2}

Decompose the total variation. Since $\hat\eps\perp\hat Y$ (from \eqref{eq:hat:orth}),
\[
  \underbrace{\lVert Y-\bar Y\mathbf 1\rVert^2}_{\mathrm{TSS}}
  = \underbrace{\lVert \hat Y-\bar Y\mathbf 1\rVert^2}_{\mathrm{ESS}}
  + \underbrace{\lVert Y-\hat Y\rVert^2}_{\mathrm{RSS}},
\]
and
\begin{equation}
  R^2 = 1-\frac{\mathrm{RSS}}{\mathrm{TSS}}\in[0,1],
  \qquad
  \bar R^2 = 1-\frac{\mathrm{RSS}/(T-n)}{\mathrm{TSS}/(T-1)} .
  \label{eq:r2}
\end{equation}
$R^2$ never decreases when a regressor is added (the column space only grows), which is
why it is useless for model comparison; $\bar R^2$ penalizes the extra parameter and can
decrease.

The joint test ``all $k$ slopes are zero'' uses
\begin{equation}
  F \;=\; \frac{(\mathrm{RSS}_{\text{restricted}}-\mathrm{RSS}_{\text{unrestricted}})/k}
              {\mathrm{RSS}_{\text{unrestricted}}/(T-n)}
  \;\sim\; F_{k,\,T-n}\ \text{under }H_0 ,
  \label{eq:ftest}
\end{equation}
which under normality follows because the numerator and denominator are independent
scaled $\chi^2$ variables. The same algebraic form, with a different restricted model,
is the Ramsey RESET test of Section~\ref{sec:nonlinear:reset} --- which is why we derive
it here.

\begin{remark}[Why we distrust $R^2$ in this project]
Regressing 16 years of Adobe log prices on three peers' log prices gives
$R^2\approx0.98$. That number is nearly meaningless: all four series are trending,
so almost any trending regressor set produces it (Section~\ref{sec:prob:spurious}).
This is why Phase~\ref{ch:phase5} compares estimators on \emph{strategy-level}
out-of-sample net Sharpe and on the out-of-sample information coefficient of $z_t$
against realized reversion returns, rather than on in-sample $R^2$.
\end{remark}

\subsection{Multicollinearity, condition number, and VIF}
\label{sec:ols:multicol}

Basket names are, by design, highly correlated --- that is what makes them a hedge.
Correlation among regressors is not a violation of any Gauss--Markov assumption; it is a
statement that the data is uninformative about certain directions of $\beta$.

\paragraph{Variance inflation factor.} Regress $X_j$ on all the other columns and let
$R^2_j$ be that regression's $R^2$. Then
\begin{equation}
  \Var(\hat\beta_j) \;=\; \frac{\sigma^{2}}{(1-R^2_j)\sum_t (x_{tj}-\bar x_j)^2}
  \;=\; \frac{\sigma^{2}\,\mathrm{VIF}_j}{\mathrm{TSS}_j},
  \qquad \mathrm{VIF}_j = \frac{1}{1-R^2_j}.
  \label{eq:vif}
\end{equation}
Derivation: partition $X=[X_j, X_{-j}]$; the $(j,j)$ element of $(X^{\top}X)^{-1}$ equals
$1/\bigl(X_j^{\top}(I-H_{-j})X_j\bigr)$ by the block-inverse formula, and
$X_j^{\top}(I-H_{-j})X_j = \mathrm{TSS}_j(1-R^2_j)$. With $R^2_j=0.98$ (typical for two
large-cap software names) $\mathrm{VIF}_j=50$: the standard error is $\sqrt{50}\approx7$
times larger than it would be with orthogonal regressors. This is precisely the pathology
ridge regression is designed to treat, and precisely why the OLS hedge weights in
Phase~\ref{ch:phase5} churn.

\paragraph{Condition number.} $\kappa(X^{\top}X)=\lambda_{\max}/\lambda_{\min}$; a common
rule of thumb is that $\kappa>30$ indicates harmful collinearity. In floating point,
relative error in the solution scales like $\kappa\cdot\varepsilon_{\text{mach}}$
($\varepsilon_{\text{mach}}\approx2.2\times10^{-16}$ for double precision), so at
$\kappa=10^{8}$ you have lost half your digits. Ridge's $\lambda I$ adds $\lambda$ to
every eigenvalue, giving $\kappa_{\text{ridge}}=(\lambda_{\max}+\lambda)/(\lambda_{\min}+\lambda)$
--- bounded by $\lambda_{\max}/\lambda$ no matter how degenerate the data are. That is a
\emph{numerical} justification for ridge that is independent of the statistical one.

\subsection{Numerical solution: Cholesky, and why we never invert}
\label{sec:ols:cholesky}

The engine must solve $(X^{\top}X+\lambda I)\beta = X^{\top}Y$ on every bar, for $n\le5$,
in well under a microsecond. The method:

\paragraph{Cholesky factorisation.} For SPD $A\in\R^{n\times n}$ there is a unique lower
triangular $L$ with positive diagonal such that $A=LL^{\top}$. Computing it
column-by-column (the Cholesky--Banachiewicz algorithm):
\begin{equation}
  L_{jj} = \sqrt{A_{jj}-\sum_{k<j}L_{jk}^2},
  \qquad
  L_{ij} = \frac{1}{L_{jj}}\Bigl(A_{ij}-\sum_{k<j}L_{ik}L_{jk}\Bigr),\quad i>j .
  \label{eq:cholesky}
\end{equation}
Cost: $\sum_{j}\bigl[(j-1) + (n-j)(j-1)\cdot 2\bigr]$ flops $= n^3/3 + O(n^2)$ --- half
the cost of Gauss--Jordan elimination, and no pivoting is needed because SPD-ness
guarantees $A_{jj}-\sum_{k<j}L_{jk}^2>0$. Then solve $Lz=b$ by forward substitution and
$L^{\top}\beta=z$ by back substitution, each $O(n^2)$.

\paragraph{Why not form $A^{-1}$?} Three reasons. (i) It costs $4n^3/3$ flops versus
$n^3/3+2n^2$. (ii) It is less accurate: the explicit inverse of an ill-conditioned
matrix has larger elementwise error, and multiplying by it compounds that. (iii) It
destroys structure: $A^{-1}$ is dense even when $A$ is banded. The engine's
\code{src/model/dense\_la.h} implements the SPD/ridge Cholesky solve and never forms an
inverse; the brute-force reference implementation used to cross-check it (to $10^{-8}$,
Phase~\ref{ch:phase4}) also solves rather than inverts.

\paragraph{Alternatives and when they matter.} QR ($X=QR$, so $\hat\beta=R^{-1}Q^{\top}Y$)
costs $2n^2T - 2n^3/3$ and is more accurate for ill-conditioned $X$ because it never
forms $X^{\top}X$ (which squares the condition number). SVD is the most robust and most
expensive. For $n\le5$ and $T=120$, forming $X^{\top}X$ costs $\kappa^2$ amplification
on a matrix whose $\kappa$ we control with $\lambda$, so Cholesky on the accumulated
normal equations is the right engineering trade --- and it is what makes the
$O(n)$-per-bar incremental update of Chapter~\ref{ch:systems} possible at all.

\subsection{Incremental updates: the algebra that makes the engine fast}
\label{sec:ols:incremental}

A rolling window of $W$ bars advances one bar at a time. Recomputing $X^{\top}X$ from
scratch each bar costs $O(Wn^2)$; maintaining it incrementally costs $O(n^2)$:
\begin{equation}
  (X^{\top}X)_{t} \;=\; (X^{\top}X)_{t-1} \;+\; x_{\text{in}}x_{\text{in}}^{\top}
  \;-\; x_{\text{out}}x_{\text{out}}^{\top},
  \qquad
  (X^{\top}y)_t = (X^{\top}y)_{t-1} + x_{\text{in}}y_{\text{in}} - x_{\text{out}}y_{\text{out}},
  \label{eq:rank1update}
\end{equation}
where $x_{\text{in}}$ is the row entering the window and $x_{\text{out}}$ the row leaving
it. Each $xx^{\top}$ is a \emph{rank-1} (outer-product) update. Two consequences:

\begin{enumerate}
  \item The per-bar work is $n^2/2$ multiply-adds for the two symmetric outer products
    ($n=4\Rightarrow 8$ flops each, 16 total) plus $n$ for the cross-product --- versus
    $Wn^2 = 120\cdot16=1920$ for a refit. That is the difference between the measured
    $\approx0.6\,\mu$s regression stage (Phase~\ref{ch:phase11}) and $\approx70\,\mu$s.
  \item \textbf{Numerical drift.} Adding and subtracting rank-1 terms for 4{,}190 bars
    accumulates rounding error: the maintained $X^{\top}X$ slowly departs from the true
    one, and can even lose positive-definiteness. This is not hypothetical --- it is the
    standard failure mode of naive incremental least squares. The engineering answer,
    implemented and tested here, is a periodic exact recomputation and a
    brute-force-versus-incremental equality assertion (Phase~\ref{ch:phase4}: agreement
    to $10^{-8}$).
\end{enumerate}

\paragraph{Sherman--Morrison--Woodbury.} If you maintained the \emph{inverse} instead, a
rank-1 update would be
\begin{equation}
  (A+uv^{\top})^{-1} = A^{-1} - \frac{A^{-1}uv^{\top}A^{-1}}{1+v^{\top}A^{-1}u},
  \label{eq:shermanmorrison}
\end{equation}
i.e.\ $O(n^2)$ instead of $O(n^3)$ --- the same idea, applied one level up. It is how a
Kalman filter can update a covariance cheaply, and it is the algebraic reason ridge
``leave-one-out'' cross-validation can be computed in closed form without refitting
(Section~\ref{sec:ridge:loo}). We do not use it in the hot path because we never maintain
inverses; we derive it because it explains why the Kalman recursion of
Chapter~\ref{ch:kalman} is cheap.

\subsection{Where least squares breaks on financial data}
\label{sec:ols:breaks}

Three failures, each with a named remedy used later in this book.

\begin{enumerate}
  \item \textbf{(G2) fails: non-stationary regressors.} Log prices are $I(1)$, so the
    CLT-based inference behind \eqref{eq:tstat} is invalid and regressions can be
    spurious. Remedy: test the residual for a unit root
    (Chapter~\ref{ch:timeseries}); if the residual is $I(0)$ the regression is
    cointegrating and inference on $\hat\beta$ is actually \emph{super-consistent}
    (converges at rate $T$, not $\sqrt T$) though non-normal.
  \item \textbf{(G3) fails: autocorrelated and heteroskedastic errors.} Daily residuals
    are strongly serially correlated (that is the mean reversion!) and strongly
    heteroskedastic (volatility clustering). Point estimates stay unbiased; standard
    errors are wrong, typically too small. Remedy: Newey--West HAC
    (Section~\ref{sec:ts:hac}), reported before/after in Phase~\ref{ch:phase6}.
  \item \textbf{$n$ large relative to $T$, or $\kappa$ large.} Variance explodes.
    Remedy: regularization (Chapter~\ref{ch:regularization}) or dimension reduction
    (PCR), and adaptive estimation when the true $\beta$ drifts
    (Chapter~\ref{ch:kalman}).
\end{enumerate}

The measured evidence for all three is in
Tables~\ref{tab:diag_stationarity}--\ref{tab:diag_coint}: Durbin--Watson $\approx0.15$
(extreme positive serial correlation), Breusch--Pagan and White $p<10^{-3}$ (reject
homoskedasticity), ADF failing to reject a unit root in the static hedge residual
($p\approx0.08$), and HAC $t$-statistics roughly $3$--$4\times$ smaller than the classic
ones ($t_{\text{CRM}}$: 29.5 $\rightarrow$ 7.3).
